\documentclass[conference]{IEEEtran}

\usepackage{cite}
\usepackage{amsmath}
\usepackage{graphicx}
\usepackage{booktabs}
\usepackage{array}
\usepackage{url}
\usepackage{listings}
\usepackage{listings}

\graphicspath{{./}{paper/}{paper/figures/}}

\begin{document}

\title{Report on the Research's Findings \\
\largeMemoLang investigates whether destructive-read semantics can provide measurable memory-management benefits for GPU matrix workloads.}

\author{
Independent Interpreter Project\\
September 2026}

\maketitle

\begin{abstract}
This paper presents MemoLang, an interpreter and programming-language design
project that investigates whether destructive-read semantics can reduce memory
and runtime overhead in matrix-oriented programs. MemoLang uses LLVM as a
backend and integrates optimized GPU matrix multiplication through NVIDIA's
CUDA ecosystem, with CUTLASS and cuBLAS considered as underlying GPU
implementations.

The primary hypothesis is not that the language can make an optimized matrix
multiplication kernel itself substantially faster. Instead, MemoLang explores
whether explicit consumption of matrix operands can reduce the lifetime of
intermediate tensor buffers, decrease temporary-buffer allocation, and lower
peak GPU memory usage. Such effects may be particularly relevant to workloads
containing chains of matrix operations or computation graphs with short-lived
intermediate tensors.

The project evaluates MemoLang against equivalent non-destructive
implementations and optimized GPU-library baselines. Measurements include
matrix multiplication time, end-to-end execution time, peak memory usage,
buffer allocation behavior, and numerical correctness. This distinction
allows improvements attributable to memory management to be separated from
the performance of the underlying GPU kernels.

The resulting evaluation is intended to determine the workloads for which
destructive-read semantics provide measurable practical benefits, as well as
the cases in which the additional language and runtime mechanisms provide
little or no advantage.
\end{abstract}

\begin{IEEEkeywords}
matrix multiplication, LLVM, destructive-read, memory efficient, Python-like syntax, C++, CUTLASS
\end{IEEEkeywords}

\section{Introduction}

Matrix multiplication is one of the fundamental operations in modern
computing. It appears in graphics, signal processing, scientific computing,
machine learning, and transformer models. Modern GPU libraries such as
cuBLAS and CUTLASS already provide highly optimized implementations of matrix
multiplication, making it difficult for a higher-level programming language
to improve the underlying arithmetic kernel without changing the generated
GPU code or algorithm.

This observation motivates the central research question of MemoLang:
rather than attempting to make the matrix multiplication kernel itself
faster, can a programming language improve the way memory surrounding matrix
operations is managed?

MemoLang is an interpreter and language design project that investigates
destructive-read semantics for matrix and tensor values. Under this model,
an operation may consume its input buffer when the value is known to have no
future consumers. The consumed storage can subsequently become eligible for
reuse by another intermediate value.

For example, in a sequence of operations

\begin{equation}
T_1 = A B,\qquad
T_2 = T_1 C,\qquad
T_3 = T_2 D,
\end{equation}

a conventional implementation may maintain several intermediate buffers at
the same time. If $A$ and $T_1$ are no longer required after their respective
operations, destructive-read semantics may allow their storage to be reused.

The expected benefit is therefore primarily a memory-management benefit,
rather than an improvement to the arithmetic throughput of an individual
matrix multiplication. If MemoLang generates the same optimized CUTLASS or
cuBLAS kernel as an equivalent implementation, the GPU kernel itself should
have approximately the same computational performance. Any end-to-end
improvement would instead result from factors such as reduced temporary
allocation, reduced buffer lifetime, or reduced memory pressure.

This distinction is important because an optimized GPU matrix multiplication
can dominate the execution time of a large workload. In such cases,
language-level optimization may have little effect on kernel execution time.
Conversely, workloads containing many intermediate tensors or relatively
small matrix operations may spend a greater fraction of their execution time
on memory management, making them more suitable for evaluating destructive
reads.

The project therefore investigates the following hypotheses:

\begin{itemize}
\item Destructive-read semantics can reduce peak memory usage for workloads
containing short-lived intermediate matrices.
\item Destructive reads can reduce the number or cost of temporary-buffer
allocations when consumed buffers can safely be reused.
\item For workloads in which memory management represents a significant
portion of total execution time, buffer reuse may improve end-to-end
performance.
\item For large matrix multiplications dominated by GPU computation,
destructive-read semantics may have little measurable effect on kernel
execution time.
\end{itemize}

This project was inspired by research into destructive-read memory systems,
the Sere programming language, the Clean programming language, and optimized
GPU matrix-multiplication libraries.

\begin{itemize}
\item Improve performance for matrix multiplication operations, supporting faster training and inference for machine learning models.
\item Lower the hardware-specification entry barrier for developers implementing machine learning algorithms and LLM-related code, by providing a high-level language that abstracts away low-level optimization details.
\item Allows compiler to effectively reuse immediate tensor buffers.
\item Eliminate run-time buffer and cache allocating overhead.
\end{itemize}


This project was inspired by the Sere programming language, the CUTLASS library, and the Clean compiler.

\section{Background}

% TODO: Motivate destructive-read semantics, explain why they help memory
% efficiency for matrix workloads, and give any relevant background on LLVM
% IR generation or CuBLAS's programming model.

% TODO: Motivate destructive-read semantics, explain why they help memory
% efficiency for matrix workloads, and give any relevant background on LLVM
% IR generation or CuBLAS's programming model.

Matrix multiplication can be written as

\begin{equation}
C_{i,j} = \sum_{k=0}^{N-1} A_{i,k}B_{k,j}
\end{equation}

Each value in the output matrix $C$ is produced by multiplying one row of $A$ with one column of $B$ and summing the products. Efficient implementations must manage data movement, memory reuse, and (in a destructive-read scheme) the lifetime of operand buffers once they are consumed.

Destructive-read is a memory read operation that inherently erases the value that it reads, which means that data is destroyed at the same time that it is accessed and read. In the case that a memory must be preserved, a write-back operation must be specified to immediately write the value back, or use external buffer/cache to restore the memory.
In modern, day, this practice is not used much in software due to the complexity of implementation. Modern day programming languages use a combination of ownership, and immutability to implement move semantics to maximize safety and minimize performance drawbacks or through other systems instead of the unsafe destructive-read. Past implementation of Linear systems ( variables must be used exactly once ). Studies show that a modified DRAM can offer lower random access time for DRAM, large performance benefits [1]. Though this is at a hardware level, we believe that applying it to software such as an Interpreter can offer the same or different emerging benefits. The research presented an optimized scheme:
" We have proposed new schemes that remove the large write-back buffer and increase performance by utilizing the cache. The data is first read from DRAM and into the cache. At this time data are not stored in DRAM, only in the cache. Data are written back to DRAM when the cache line is replaced."
Judging by the emerged performance benefit on the modified DRAM, we believe a software-level implementation could optimize buffer and cache allocating as well as an increase in performance.
Though the domain of destructive-read exist largely in hardware and circuit architecture, we believe that implementing it for software, specifically one where optimizing performance is at foremost concern, could offer universal benefits for both software developers and users.
	% TODO: Add background on destructive read: what it means for a value to be
	% "consumed" once read, how this differs from ordinary immutable/borrowed
	% semantics, and what memory-safety guarantees (if any) the language provides. BUFFER AS BIG AS DRAM BANK SIZE
	%" We have proposed new schemesthat remove the large write-back buffer and increaseperformance by utilizing the cache [5]. The data is firstread from DRAM and into the cache. At this time dataare not stored in DRAM, only in the cache. Data are written back to DRAM when the cache line is replaced."

\section{Project Goal}

The goal of MemoLang is to determine whether language-level destructive-read
semantics can provide measurable memory-management benefits for
matrix-oriented GPU programs while retaining the performance of optimized
matrix multiplication kernels.

The project will be implemented and evaluated through the following goals:

\begin{itemize}
\item Define MemoLang's core Python-like syntax and destructive-read
semantics.
\item Implement a lexer, parser, AST, and LLVM-based code-generation
pipeline.
\item Generate LLVM IR for arithmetic and matrix operations.
\item Integrate optimized GPU matrix multiplication through CUTLASS and/or
cuBLAS.
\item Implement explicit ownership and consumption tracking for matrix
operands.
\item Determine when consumed matrix buffers can safely be reused.
\item Verify generated results against an independent numerical reference.
\item Measure peak GPU memory usage and temporary-buffer allocation behavior.
\item Measure both GPU kernel execution time and end-to-end program runtime.
\item Compare destructive-read execution against an equivalent
non-destructive implementation.
\item Compare MemoLang against optimized GPU-library baselines where
appropriate.
\item Preserve benchmark logs, correctness results, compiler output, and
build artifacts to make the experimental results reproducible.
\end{itemize}

The evaluation will not assume that destructive reads provide a performance
improvement. A reduction in peak memory usage without a reduction in runtime
will be treated as a meaningful result, while an observed runtime improvement
will be analyzed to determine whether it originates from reduced memory
management overhead rather than from changes to the underlying matrix
multiplication kernel.

\subsection{Destructive-Read Semantics}

A destructive read treats a value as consumed when an operation has exclusive
use of that value. After consumption, the original value is no longer
available to subsequent operations unless the programmer explicitly creates a
copy or otherwise preserves it.

For matrix operations, this provides an opportunity for storage reuse. For
example,

\begin{lstlisting}
C = A @ B
D = C @ W
\end{lstlisting}

may permit the storage associated with $A$ or $C$ to become eligible for
reuse once the compiler establishes that the corresponding value has no
remaining consumers.

The purpose of this mechanism is not to alter the mathematical computation
performed by matrix multiplication. Instead, it changes the lifetime of
buffers associated with the computation. If a buffer can be reused without
copying its contents, the implementation may reduce peak memory consumption
and allocation activity.

This optimization is inherently workload-dependent. A matrix with multiple
future consumers cannot safely be destroyed after its first use. In such
cases, the implementation may require a copy, borrow, or other mechanism to
preserve the value. Consequently, destructive-read semantics can introduce a
tradeoff between memory efficiency and the flexibility to reuse values.

The language therefore treats destructive reads as an explicit ownership
mechanism rather than assuming that consuming every operand will necessarily
improve performance.

% TODO: Describe MemoLang's syntax and semantics. Suggested subsections:
% - Surface syntax and Python-like conventions
% - Type system (if any) for matrices/tensors
% - Destructive-read semantics: how a read consumes a value, what happens on
%   re-use, error/ownership model
% - How CUDA's programming model influenced the design

\subsection{Syntax Overview}
% TODO: Example MemoLang snippet(s) illustrating matrix declarations and
% multiplication, e.g.:
%
% \begin{lstlisting}
% let A = matrix(4, 4)
% let B = matrix(4, 4)
% let C = A @ B   # destructive read of A and B
% \end{lstlisting}

\subsection{Destructive-Read Semantics}
% TODO: Formal or informal description of destructive read: once a matrix
% value is read (e.g., consumed by @), it can no longer be used unless
% explicitly cloned. Explain how this reduces peak memory usage.

\section{Compiler and Runtime Architecture}

% TODO: Describe the pipeline: source -> lexer -> parser -> ExprAST ->
% LLVM IR -> JIT/AOT compilation -> CuBLAS call for MatMul kernels.
% A table of main source files (mirroring Table~\ref{tab:modules} in the
% original hardware paper, but for this project) belongs here, e.g.:
\section{Maths and Algorithms}

% TODO: Describe MemoLang's syntax and semantics. Suggested subsections:
% - Surface syntax and Python-like conventions
% - Type system (if any) for matrices/tensors
% - Destructive-read semantics: how a read consumes a value, what happens on
%   re-use, error/ownership model
% - How CUDA's programming model influenced the design

\subsection{Math behind FP32 Matrix Multiplication}
% TODO: Example MemoLang snippet(s) illustrating matrix declarations and
% multiplication, e.g.:
%
% \begin{lstlisting}
% let A = matrix(4, 4)
% let B = matrix(4, 4)
% let C = A @ B   # destructive read of A and B
% \end{lstlisting}

\subsection{The "special scheme" of destructive-read semantics}
% TODO: Formal or informal description of destructive read: once a matrix
% value is read (e.g., consumed by @), it can no longer be used unless
% explicitly cloned. Explain how this reduces peak memory usage.

\section{CUTLASS Integration}

% TODO: Describe the pipeline: source -> lexer -> parser -> ExprAST ->
% LLVM IR -> JIT/AOT compilation -> CUTLASS call for MatMul kernels.
% A table of main source files (mirroring Table~\ref{tab:modules} in the
% original hardware paper, but for this project) belongs here, e.g.:

\begin{table}[h]
\centering
\caption{Main Source Files}
\caption{Main Source Files}
\label{tab:modules}
\begin{tabular}{ll}
\toprule
File & Purpose \\
File & Purpose \\
\midrule
\texttt{Lexer.h/.cpp} & Tokenizes MemoLang source \\
\texttt{Parser.h/.cpp} & Builds the AST \\
\texttt{ExprAST.h} & AST node definitions \\
\texttt{CodeGen.cpp} & LLVM IR generation \\
\texttt{CuBLASRuntime.cpp} & CuBLAS MatMul bindings \\
\texttt{main.cpp} & Driver / REPL \\
\texttt{Lexer.h/.cpp} & Tokenizes MemoLang source \\
\texttt{Parser.h/.cpp} & Builds the AST \\
\texttt{ExprAST.h} & AST node definitions \\
\texttt{CodeGen.cpp} & LLVM IR generation \\
\texttt{CuBLASRuntime.cpp} & CuBLAS MatMul bindings \\
\texttt{main.cpp} & Driver / REPL \\
\bottomrule
\end{tabular}
\end{table}

% TODO: Describe how MatMul expressions lower to LLVM IR, and how/when
% control transfers to CUTLASS (e.g., via an external function call inserted
% during codegen). Discuss the destructive-read buffer handoff to CUTLASS.

\section{Verification}

% TODO: Describe how correctness was tested (e.g., comparing MemoLang
% MatMul output against NumPy or cuBLAS reference results for various
% matrix sizes). Replace with a results table once available, e.g.:
% TODO: Describe how correctness was tested (e.g., comparing MemoLang
% MatMul output against NumPy or cuBLAS reference results for various
% matrix sizes). Replace with a results table once available, e.g.:

\begin{table}[h]
\centering
\caption{Correctness Test Results (Template)}
\label{tab:verification}
\caption{Correctness Test Results (Template)}
\label{tab:verification}
\begin{tabular}{ll}
\toprule
Test Case & Result \\
Test Case & Result \\
\midrule
$N \times N$ MatMul vs. reference & TBD \\
Destructive-read re-use error detection & TBD \\
Edge cases (empty, 1x1, non-square) & TBD \\
$N \times N$ MatMul vs. reference & TBD \\
Destructive-read re-use error detection & TBD \\
Edge cases (empty, 1x1, non-square) & TBD \\
\bottomrule
\end{tabular}
\end{table}

\section{Benchmarks}

The benchmark evaluation is designed to distinguish improvements in memory
management from improvements in matrix multiplication itself.

Three primary implementations will be compared:

\begin{enumerate}
\item An optimized GPU-library implementation used as a kernel-performance
reference.
\item MemoLang using conventional non-destructive value semantics.
\item MemoLang using destructive-read semantics.
\end{enumerate}

Where appropriate, a simple C++ or CUDA implementation will also be included
as a baseline.

The most important comparison is between the two MemoLang configurations.
Both configurations perform the same mathematical operations and use the
same underlying GPU matrix multiplication implementation. This comparison
isolates the effect of destructive-read semantics more effectively than a
comparison against an unrelated implementation.

The benchmark suite will include isolated matrix multiplications as well as
chains of matrix operations. Matrix dimensions will include small, medium,
large, and non-square matrices. This allows the evaluation to determine
whether the effect of destructive reads depends on workload size.

The following measurements will be collected:

\begin{itemize}
\item GPU kernel execution time.
\item End-to-end execution time.
\item Peak GPU memory consumption.
\item Number of temporary-buffer allocations.
\item Number of buffer reuse events.
\item Host-to-device and device-to-host transfer time where applicable.
\item Compiler and runtime overhead.
\item Numerical error relative to a reference implementation.
\end{itemize}

GPU measurements will use multiple repetitions and appropriate warm-up
iterations. CUDA initialization and compilation overhead will be separated
from steady-state execution measurements where possible.

The benchmark will report both absolute measurements and relative changes.
No performance improvement will be attributed to destructive-read semantics
unless the corresponding measurements support that interpretation.

% TODO: Report runtime and peak-memory measurements comparing:
% (a) naive/non-destructive baseline, (b) MemoLang with destructive read,
% (c) raw CuBLAS call for reference. A results table/figure goes here.
\section{Benchmarks}

% TODO: Report runtime and peak-memory measurements comparing:
% (a) naive/non-destructive baseline, (b) MemoLang with destructive read,
% (c) raw CuBLAS call for reference. A results table/figure goes here.

\section{Evidence and Reproducibility}

% TODO: List where logs, benchmark results, and build artifacts are saved
% (mirroring the "Evidence and Reproducibility" approach in hardware
% projects) so claims in this paper are traceable to saved output.
% TODO: List where logs, benchmark results, and build artifacts are saved
% (mirroring the "Evidence and Reproducibility" approach in hardware
% projects) so claims in this paper are traceable to saved output.

\subsection{Real-World Benefits}

The primary potential real-world benefit of MemoLang is improved memory
efficiency rather than faster individual matrix multiplication kernels.

Modern GPU libraries already provide highly optimized matrix multiplication
kernels. If MemoLang ultimately invokes the same kernel with equivalent
inputs, there is little reason to expect destructive-read semantics alone to
significantly increase the arithmetic throughput of that kernel.

The potential advantage instead occurs when multiple operations create
intermediate tensors. Consider a computation such as

\begin{equation}
T_1 = AB,\qquad
T_2 = T_1C,\qquad
T_3 = T_2D.
\end{equation}

A conventional implementation may need to maintain several intermediate
buffers simultaneously. If each intermediate has only one consumer, a
destructive-read implementation can potentially reuse storage after that
consumer has completed.

This can produce practical benefits even if individual matrix multiplication
kernels execute at exactly the same speed. Lower peak memory consumption can
allow larger batch sizes, larger sequence lengths, or larger models to fit
within a fixed GPU memory capacity. Reducing allocation and deallocation
activity may also reduce runtime overhead in workloads containing many small
operations.

However, these benefits are not universal. Machine-learning workloads often
contain residual connections, branching computation graphs, parameter reuse,
concatenation, normalization, and other operations in which values have
multiple consumers. Such workloads may require values to remain alive or may
require explicit copies, potentially reducing or eliminating the memory
advantage.

The practical value of MemoLang should therefore be evaluated according to
the workload rather than expressed as a universal performance improvement.

\subsection{Expected Relationship Between Language and MatMul Performance}

The expected execution time can be approximately decomposed as

\begin{equation}
T_{\mathrm{total}} =
T_{\mathrm{language}} +
T_{\mathrm{memory}} +
T_{\mathrm{launch}} +
T_{\mathrm{kernel}}.
\end{equation}

Destructive-read semantics primarily target
$T_{\mathrm{memory}}$. They are not expected to substantially reduce
$T_{\mathrm{kernel}}$ when the same optimized GPU kernel is used.

For a large matrix multiplication, $T_{\mathrm{kernel}}$ may dominate the
total execution time, resulting in little observable end-to-end improvement.
For workloads containing many intermediate tensors or small matrix
operations, however, $T_{\mathrm{memory}}$ and other runtime costs may become
more significant.

Consequently, a successful result does not require MemoLang to make an
individual matrix multiplication faster. A reduction in peak memory usage,
allocation activity, or end-to-end overhead can constitute a meaningful
result even when kernel execution time remains approximately unchanged.
\section{Threats to Validity}

The proposed benefits of destructive-read semantics should not be interpreted
as guaranteed improvements for all matrix workloads.

First, optimized GPU matrix multiplication libraries already perform
substantial low-level optimization. If MemoLang ultimately invokes the same
kernel with the same input dimensions and precision, destructive-read
semantics should not be expected to substantially change the arithmetic
throughput of that kernel.

Second, memory savings do not necessarily imply runtime savings. A program
may use less memory while spending approximately the same amount of time
performing the underlying matrix operations. Conversely, additional ownership
tracking or copying could make some workloads slower.

Third, benchmark results depend on GPU architecture, CUDA version, library
version, matrix dimensions, precision, synchronization behavior, allocation
strategy, and data-transfer policy. These variables must therefore be
reported and controlled when comparing implementations.

Finally, improvements measured on synthetic matrix chains may not directly
translate to complete machine-learning systems. Real computation graphs
contain values with different numbers of consumers and different lifetime
requirements. Results should therefore be interpreted in terms of the
specific workloads tested rather than generalized to all GPU applications.

\section{Future Work}

% TODO: Outline next steps -- broader operator support, multi-GPU,
% integration with autodiff, formal semantics for destructive read, etc.
% TODO: Outline next steps -- broader operator support, multi-GPU,
% integration with autodiff, formal semantics for destructive read, etc.

\section{Conclusion}

% TODO: Summarize what was built and verified, and what remains as future
% work, tying claims back to saved evidence as in Section~\ref{tab:evidence}
% above (adjust reference once that section is filled in).
% TODO: Summarize what was built and verified, and what remains as future
% work, tying claims back to saved evidence as in Section~\ref{tab:evidence}
% above (adjust reference once that section is filled in).

\begin{thebibliography}{1}
\bibitem{DybdahlKGN06}
Haakon Dybdahl, Per Gunnar Kjeldsberg, Marius Grannæs, and Lasse Natvig,
``Destructive-read in embedded DRAM, impact on power consumption,''
\emph{Journal of Embedded Computing}, vol.~2, no.~1, pp.~83--96, 2006.
\url{http://dblp.uni-trier.de/rec/bibtex/journals/jec/DybdahlKGN06}
\bibitem{cublas}
NVIDIA, ``cuBLAS Library.''

\bibitem{llvm}
LLVM Project, ``The LLVM Compiler Infrastructure.''
\bibitem{llvm}
LLVM Project, ``The LLVM Compiler Infrastructure.''

\bibitem{clean}
Clean Language Team, ``The Clean Programming Language.''
\bibitem{clean}
Clean Language Team, ``The Clean Programming Language.''

\bibitem{sere}
Sere Language Project, ``Sere Programming Language.''
\bibitem{sere}
Sere Language Project, ``Sere Programming Language.''

\bibitem{memolang}
Independent Interpreter Project, ``MemoLang: source and benchmark evidence,'' project repository, September 2026.
\bibitem{memolang}
Independent Interpreter Project, ``MemoLang: source and benchmark evidence,'' project repository, September 2026.

\end{thebibliography}

\end{document}